% Original PDF page 14; hash in recovery-map.json.
\recoverypage{14}
\section*{C  Modal semantics and adversarial minimal pairs}
% Editable layout transcription; cell boundaries not asserted.
\begin{center}\resizebox{\linewidth}{!}{\begin{tabular}{@{}l@{}}
\hspace*{16\fontdimen2\font}Table 3: Minimal distinctions to preserve across source, IR, and backend. \\
 \\
Distinction\hspace*{17\fontdimen2\font}Counterexample or ambiguity\hspace*{17\fontdimen2\font}Correct boundary \\
Obligation vs fact\hspace*{6\fontdimen2\font}A required event has not happened. Ideal\hspace*{8\fontdimen2\font}Do not infer p from O(p). \\
\hspace*{24\fontdimen2\font}worlds satisfy p while the actual world \\
\hspace*{24\fontdimen2\font}does not. \\
Permission vs occur- An action is allowed but never taken.\hspace*{14\fontdimen2\font}Permission is not an observa- \\
rence\hspace*{67\fontdimen2\font}tion. \\
Prohibition vs negation A forbidden action can occur and violate\hspace*{8\fontdimen2\font}Keep normative force distinct \\
\hspace*{24\fontdimen2\font}the norm.\hspace*{39\fontdimen2\font}from actual falsity. \\
Knowledge vs belief\hspace*{5\fontdimen2\font}An agent believes an incorrect claim.\hspace*{11\fontdimen2\font}Factivity requires the knowl- \\
\hspace*{72\fontdimen2\font}edge profile and appropriate \\
\hspace*{72\fontdimen2\font}frame condition. \\
Always vs now\hspace*{15\fontdimen2\font}p now, not p at a later step.\hspace*{15\fontdimen2\font}Preserve the trace interval, not \\
\hspace*{72\fontdimen2\font}just p. \\
Until vs conjunction\hspace*{8\fontdimen2\font}q holds now while p does not. Strong\hspace*{8\fontdimen2\font}Preserve temporal operator se- \\
\hspace*{28\fontdimen2\font}until holds, conjunction does not.\hspace*{10\fontdimen2\font}mantics. \\
Quantifier order\hspace*{12\fontdimen2\font}Each user has some key, but no key be-\hspace*{6\fontdimen2\font}Preserve binder order and \\
\hspace*{28\fontdimen2\font}longs to every user.\hspace*{24\fontdimen2\font}identity. \\
Only if vs if\hspace*{15\fontdimen2\font}``Write only if approved'' differs from ``if\hspace*{3\fontdimen2\font}Check implication direction \\
\hspace*{28\fontdimen2\font}approved then write.''\hspace*{23\fontdimen2\font}independently. \\
Exception vs conjunc-\hspace*{7\fontdimen2\font}``Must audit unless exempt'' is not un-\hspace*{7\fontdimen2\font}Preserve exception scope and \\
tion\hspace*{24\fontdimen2\font}conditional auditing or an audit of only\hspace*{4\fontdimen2\font}priority. \\
\hspace*{28\fontdimen2\font}exempt cases. \\
Observed absence vs\hspace*{9\fontdimen2\font}A partial log has no audit event.\hspace*{11\fontdimen2\font}Unknown\hspace*{8\fontdimen2\font}unless\hspace*{5\fontdimen2\font}cap- \\
negation\hspace*{64\fontdimen2\font}ture/deadline completeness is \\
\hspace*{72\fontdimen2\font}established. \\
Same spelling vs iden-\hspace*{6\fontdimen2\font}A log's ``admin'' and a policy's ``admin''\hspace*{6\fontdimen2\font}Reject an ungrounded entity \\
tity\hspace*{24\fontdimen2\font}belong to different tenants.\hspace*{16\fontdimen2\font}join. \\
Round-trip vs fidelity\hspace*{6\fontdimen2\font}Both compiler and realizer drop the same\hspace*{4\fontdimen2\font}Add independent semantic an- \\
\hspace*{28\fontdimen2\font}exception.\hspace*{34\fontdimen2\font}notation and negative cases.
\end{tabular}}\end{center}

\noindent
The core registry names additional profiles, including dynamic, conditional normative, frame, and\linebreak[4]
hybrid views. Before evaluating such a view, state its actual interpretation. An action modality may\linebreak[4]
use a transition relation; an event-calculus fluent uses time and event axioms; a frame-based ontology\linebreak[4]
may use inheritance and method relations. A cognitive operator may be hyperintensional rather than\linebreak[4]
an ordinary Kripke box. These are design choices to bind, not silently infer from an operator name.\par

\section*{D  Executable reference semantics}
\noindent
The included evaluation/reference\_semantics.py uses only the Python standard library. It\linebreak[4]
checks the paper's finite examples, not the project modules. The expected-result file is generated by\linebreak[4]
the same script and is therefore a reproducibility artifact, not an independent assurance certificate.\linebreak[4]
No timed performance or learned accuracy claim is derived from it.\par

\noindent
python evaluation/reference\_semantics.py\linebreak[4]
--output evaluation/reference\_results.json\par

\noindent
For the illustrative normal deontic check only, O is a universal modality and P its existential dual,\linebreak[4]
P p = \ensuremath{\neg}O\ensuremath{\neg}p; seriality supports the tested implication. This does not equate strong explicit permission\linebreak[4]
with weak dual permission in every normative logic. The suite enumerates all 16 relations on two\linebreak[4]
worlds and the relevant valuations for normal modal distribution; reflexive relations for factivity;\linebreak[4]
serial relations for the D implication; all length-three Boolean traces for strong until; and finite\linebreak[4]
request/trace relations for the worked example. Additional explicit witnesses distinguish obligations\par

